\cotefile{62}
\scanpage{20}
% archive sheet number: 19
\mspage{- 20 -}\hfill n$^{\mathrm{o}}$ 276

\handwritten{\S{}}2. - \uline{Anneaux} $\mathrm{m}$-\uline{adiques noethériens}.

Dans tout ce \handwritten{\S{}} on désigne par $A$ un anneau \uline{commutatif noethérien} et par $\mathrm{m}$ un idéal de $A$ ; sauf mention explicite du contraire on munira $A$, et plus généralement tous les $A$-modules considérés, de la topologie $\mathrm{m}$-adique.

1. - \uline{Restriction de la topologie} $\mathrm{m}$-\uline{adique à un sous-module}.

\uline{Théorème} 1 - \uline{Soient} $E$ \uline{un} $A$-\uline{module de type fini et} $F$ \uline{un sous-module de} $E$. \uline{La topologie} $\mathrm{m}$-\uline{adique du} $A$-\uline{module} $F$ \uline{est induite par la topologie} $\mathrm{m}$-\uline{adique de} \unsure{-|} $E$.

Le théorème 1 sera une conséquence du résultat plus précis que voici :

\uline{Lemme d'Artin-Rees} $\handwritten{\unsure{\llcorner|-}}$ \uline{Soient} $A$ \uline{un anneau commutatif noethérien}, $\mathrm{m}$ \uline{un idéal de} $A$, $E$ \uline{un} $A$-\uline{module de type fini et} $F$ \uline{un sous-module de} $E$. \uline{Pour tout entier} $k$ \uline{suffisamment grand on a}
\[
\mathrm{m}^{p}E\handwritten{\cap} F=\mathrm{m}^{p-k}(\mathrm{m}^{k}E\handwritten{\cap} F)\qquad \text{\uline{pour tout}}~p\handwritten{\geqslant} k.
\]

Il suffit évidemment de prouver l'inclusion \unsure{-|}
\[
\mathrm{m}^{p}E\handwritten{\cap} F\handwritten{\subset}\mathrm{m}^{p-k}(\mathrm{m}^{k}E\handwritten{\cap} F)\quad,
\]
puisque l'inclusion opposée est triviale.

Pour cela considérons dans l'anneau de polynômes $A\handwritten{[}X\handwritten{]}$ le sous-anneau $A'$ formé des polynômes $\handwritten{\sum} a_p X^p$ tels que $a_p\handwritten{\in}\mathrm{m}^p$ pour tout $p$. Il est clair qu'en tant qu'algèbre sur $A$, $A'$ est engendré par les $a_i X$ où $a_i$ est un système de générateurs de $\mathrm{m}$ ; par suite $A'$ est un anneau noethérien, par ailleurs gradué.

Formons de même le $A\handwritten{[}X\handwritten{]}$-module
\[
E\handwritten{[}X\handwritten{]}=A\handwritten{[}X\handwritten{]}\underset{A}{\handwritten{\otimes}}E
\]
des polynômes à coefficients dans $E$ ; en tant que $A'$-module, $E\handwritten{[}X\handwritten{]}$ contient le sous-module $E'$ formé des polynomes $\handwritten{\sum} e_p X^p$ tels que $e_p\handwritten{\in}\mathrm{m}^p E$

